\par\SourceLeaf{156}
Zu der Rechnung habe ich die höchst schätzbaren Mortalitätstafeln angewandt, welche \textsc{Brune} aus den bei der preussischen Witwenkasse an 31500 Ehepaaren gemachten Erfahrungen abgeleitet hat. Nur für das Absterben der Männer im höchsten Lebensalter sind diese Tafeln mangelhaft, da die Registratur der Witwenkasse dazu keine hinreichende Daten enthielt, und ich habe daher vorgezogen, das Absterben der Männer über 80 Jahr nach denselben Verhältnissen zu rechnen, welche die Tafeln für das weibliche Geschlecht angeben. Zur Berechnung der Modificationen, welche der Werth der Witwenpensionen durch Berücksichtigung der minderjährigen Kinder, wo solche vorhanden sind, erleidet, habe ich für letztere die Mortalitätstafeln von \textsc{Deparcieux} angewandt. Übrigens hat die Wahl der Tafeln, sowohl für die Mortalität der Kinder, als für die der Männer im höchsten Lebensalter auf die Rechnungsresultate nur sehr geringen Einfluss.

Alle Resultate sind wesentlich abhängig von dem dabei zum Grunde gelegten Zinsfuss; und ich habe daher, um hier nichts zu wünschen übrig zu lassen, sämmtliche Rechnungen sowohl nach dem Zinsfuss von 4 Procent, als nach dem von \(3\tfrac{1}{2}\) Procent durchgeführt, obgleich dadurch die Arbeit gerade verdoppelt wurde. Die Münzen sind immer als Gold zu verstehen. Die Rechnung ist durchgehends auf Brüche des Thalers genau geführt, diese Brüche aber sind in gegenwärtiger Abschrift weggelassen. Daraus wird hin und wieder ein Unterschied von einer oder ein paar Einheiten in den Summationen erscheinen können, welchen geringfügigen Umstand ich hier bloss deswegen bemerke, damit nicht jemand, der etwa die eine oder die andere der Summationen nachrechnet, solche Unterschiede für Zeichen von ungenauer Rechnung halte, da sie vielmehr gerade umgekehrt die Folge der in der Rechnung beobachteten grössern Schärfe sind.

\Needspace{6\baselineskip}\pdfbookmark[3]{I. Berechnung des Geldwerths der Witwenpensionen}{witwen-II-2-I}
\par\medskip\noindent I.\quad Berechnung des auf den 1. October 1845 reducirten Geldwerths der Witwenpensionen.\par\smallskip
Die Zahl aller seit Errichtung der Witwenkasse bis jetzt eingetretenen Pensionirungen ist 68, und ich habe dieselben, um in meiner Arbeit die leichteste Übersichtlichkeit zu gewinnen, mit fortlaufender Numerirung bezeichnet. Die Ansätze in dem folgenden tabellarischen Abrisse drücken die Summen aus, welche den einzelnen Witwen ausgezahlt werden müssten, wenn sie für ihre Ansprüche abgefunden werden sollten, oder umgekehrt, die Summen, welche die Witwen einzahlen müssten, wenn sie eine solche Berechtigung, wie ihnen jetzt zusteht, sich erst erkaufen wollten. Der Rechnung liegt der jetzige Pensionssatz von 250 Rthl. zum Grunde, mit Ausnahme der Nr. 58, wo er 200 Rthl. beträgt. Die Witwen 50. 53. 58. 64. 68 haben Kinder unter 20 Jahren, die bei der Rechnung genau berücksichtigt sind, eben so wie der Umstand, dass die Pensionen immer noch für den vollen Monat, in welchem ein Abgang statt findet, ausgezahlt werden.

\begin{center}
\footnotesize\setlength{\tabcolsep}{3pt}
\begin{tabular}{r|c|r|r||r|c|r|r}
\multicolumn{2}{c|}{}&\multicolumn{2}{c||}{\shortstack{Jetziger Werth\\der Pension.}}&\multicolumn{2}{c|}{}&\multicolumn{2}{c}{\shortstack{Jetziger Werth\\der Pensionen.}}\\\cline{3-4}\cline{7-8}
Nro.& &\(3\tfrac{1}{2}\) Proc.&4 Proc.&Nro.& &\(3\tfrac{1}{2}\) Proc.&4 Proc.\\\hline
29&L.&957 Thl.&941 Thl.&58&G.&3172 Thl.&3000 Thl.\\
31&S.&1756 ,, &1708 ,, &59&W.&2850 ,, &2732 ,, \\
41&W.&1122 ,, &1101 ,, &60&S.&2606 ,, &2507 ,, \\
44&F.&2543 ,, &2448 ,, &62&B.&1648 ,, &1605 ,, \\
50&H.&3987 ,, &3759 ,, &63&G.&2113 ,, &2045 ,, \\
51&S.&735 ,, &725 ,, &64&M.&4170 ,, &3932 ,, \\
52&P.&2352 ,, &2270 ,, &65&H.&3026 ,, &2894 ,, \\
53&M.&2789 ,, &2676 ,, &66&H.&1916 ,, &1859 ,, \\
55&H.&2813 ,, &2698 ,, &68&M.&4175 ,, &3936 ,, \\\cline{5-8}
57&S.&1798 ,, &1748 ,, &\multicolumn{2}{c}{Summa}&46529 Thl.&44582 Thl.
\end{tabular}
\end{center}
